\section{Method}
\label{sec:method}

Our goal is to improve how weak temporal supervision is used during surgical VLP without replacing the underlying CLIP-style model. The key assumption we relax is simple: a weakly paired interval should not be treated as a uniformly reliable positive clip. Instead, we treat it as a noisy temporal anchor whose supporting evidence may lie in only part of the interval and may be slightly shifted around the annotated boundary. SurgAlign addresses this issue through training-time temporal denoising.

\subsection{Problem Formulation}
\label{sec:method_setup}

Each training sample is written as
\begin{equation}
(V, [t_s, t_e], y, \ell),
\end{equation}
where $V$ is a surgical video, $[t_s, t_e]$ is the weak temporal interval, $y$ is the paired text, and $\ell \in \{f, m, c\}$ denotes the fine, mid, or coarse annotation level. In a standard CLIP-style baseline, $M$ frames are sampled from $[t_s, t_e]$, encoded by a visual backbone $E_v$, aggregated into one clip representation, and aligned with a text representation produced by a text encoder $E_t$. If $A_v(\cdot)$ is the temporal aggregation head and $P_v(\cdot), P_t(\cdot)$ are the projection heads, the baseline computes
\begin{equation}
v_{\mathrm{glob}} = \operatorname{norm}\!\left(P_v(A_v(E_v(c)))\right), \qquad
t = \operatorname{norm}\!\left(P_t(E_t(y))\right),
\end{equation}
and optimizes a symmetric contrastive loss.

This formulation is effective when the sampled clip is already well aligned with the text. Our setting is weaker: the pair may be globally correct while the truly supportive visual evidence occupies only a subset of frames. SurgAlign keeps the same backbone and contrastive formulation, but changes how the training clip is constructed and weighted.

\subsection{Overview}
\label{sec:method_overview}

The overall pipeline is illustrated in the framework diagram. Starting from a weak temporal interval, SurgAlign first expands the interval into a bounded neighborhood. It then samples candidate frames from this expanded window and extracts frame-level features. Instead of uniformly pooling these frames, SurgAlign computes text-conditioned frame scores, converts them into soft weights, and builds a selected local visual view. The model is trained with two complementary branches:
\begin{enumerate}[leftmargin=*]
    \item a global branch that preserves stable interval-level video-text alignment; and
    \item a selected branch that emphasizes text-relevant evidence inside the expanded temporal neighborhood.
\end{enumerate}
In hierarchical training, SurgAlign also encourages adjacent levels from the same video to produce compatible selected views. All of these mechanisms are training-only. Test-time inference remains standard CLIP-style video/text encoding.

\subsection{Expanded Temporal Evidence}
\label{sec:method_expand}

Weak narration timestamps often miss the most informative visual moment by a small temporal offset. We therefore replace the original interval with a bounded expanded interval
\begin{equation}
[t_s', t_e'] = \operatorname{Expand}([t_s, t_e]; \rho),
\end{equation}
where $\rho \geq 1$ is an expansion ratio and the interval is clipped to the valid video duration. We then sample a multi-frame candidate clip
\begin{equation}
c_{\mathrm{exp}} = \{x_1', \ldots, x_M'\}
\end{equation}
from $[t_s', t_e']$.

This step is deliberately modest. SurgAlign does not attempt full temporal grounding. The purpose of expansion is to give the model a small search space around a noisy timestamp so that the true supporting frames remain available during training even when the annotated boundary is slightly off.

\subsection{Multi-Level Annotation-Aware Evidence Selection}
\label{sec:method_selection}

Let the visual encoder produce frame tokens
\begin{equation}
\{h_1, \ldots, h_M\} = E_v(c_{\mathrm{exp}}),
\end{equation}
and let the text encoder produce token-level hidden states and a pooled global text representation
\begin{equation}
\{w_1, \ldots, w_L\}, \quad t = E_t(y).
\end{equation}
The global text embedding $t$ is still used by the contrastive objective. The selected branch, however, uses the token-level states to construct text-conditioned temporal evidence inside the expanded window. Our key design choice is to treat the multi-level annotation label $\ell$ as a prior on \emph{how} this evidence should be selected. Fine-, mid-, and coarse-level captions differ not only in temporal scope but also in internal semantic complexity. Fine captions are usually localized and dominated by one salient cue, whereas mid/coarse captions more often summarize multiple actions, instruments, or procedural states within one interval. We therefore treat mid/coarse captions as \emph{compositional semantic structures} rather than as longer monolithic queries.

We encode this prior through a level-conditioned evidence-selection rule. Given level $\ell$, the selected branch uses a pair of level-dependent parameters,
\begin{equation}
K_\ell, \qquad \tau_{\mathrm{sel}}^{(\ell)},
\end{equation}
where $K_\ell$ controls how many latent semantic cues may be extracted from the text and $\tau_{\mathrm{sel}}^{(\ell)}$ controls how sharply the temporal evidence is concentrated. Concretely, we use
\begin{equation}
K_\ell =
\begin{cases}
K_f, & \ell = f,\\
K_m, & \ell = m,\\
K_c, & \ell = c,
\end{cases}
\qquad
\tau_{\mathrm{sel}}^{(\ell)} =
\begin{cases}
\tau_f, & \ell = f,\\
\tau_m, & \ell = m,\\
\tau_c, & \ell = c,
\end{cases}
\end{equation}
with $K_f \leq K_m \leq K_c$ and $\tau_f < \tau_m < \tau_c$. This means that fine captions use fewer queries and sharper temporal selection, while broader captions are allowed to mine multiple latent semantic cues and distribute their mass over a wider temporal neighborhood.

To keep the training objective simple, we do not explicitly decompose the sentence into supervised sub-phrases. Instead, we introduce a small bank of learnable latent text-query slots
\begin{equation}
\{q_1, \ldots, q_{K_{\max}}\},
\end{equation}
which are shared across levels and softly attend to the token sequence to extract multiple latent semantic cues.

For the $k$-th latent query, we compute token attention
\begin{equation}
\beta_{kr} =
\operatorname{softmax}_{r}
\left(
\frac{q_k^\top w_r}{\sqrt{d}}
\right),
\end{equation}
and obtain a query-specific text cue
\begin{equation}
u_k = \sum_{r=1}^{L} \beta_{kr} w_r .
\end{equation}
These latent cues are then matched against the frame tokens after lightweight projections:
\begin{equation}
\tilde{u}_k = \operatorname{norm}(Q_t(u_k)), \qquad
\tilde{h}_j = \operatorname{norm}(Q_v(h_j)).
\end{equation}
Their similarity defines query-conditioned frame scores
\begin{equation}
s_{kj} = \tilde{u}_k^\top \tilde{h}_j .
\end{equation}

Given the annotation level $\ell \in \{f,m,c\}$, we activate only the first $K_\ell$ latent queries and normalize the scores over time with the corresponding selection temperature:
\begin{equation}
\alpha_{kj} =
\operatorname{softmax}_{j}
\left(
\frac{s_{kj}}{\tau_{\mathrm{sel}}^{(\ell)}}
\right), \qquad k = 1, \ldots, K_\ell .
\end{equation}
Each latent query therefore selects its own temporally weighted evidence from the expanded window. We aggregate the query-specific evidence as
\begin{equation}
z_k = \sum_{j=1}^{M} \alpha_{kj} h_j,
\end{equation}
and fuse them into a single selected visual view
\begin{equation}
z_{\mathrm{sel}} = \frac{1}{K_\ell}\sum_{k=1}^{K_\ell} z_k, \qquad
v_{\mathrm{sel}} = \operatorname{norm}(P_v(z_{\mathrm{sel}})).
\end{equation}

This selected branch should be interpreted as a denoised auxiliary view of the same weak pair, not as a new inference-time matching module. The global branch still provides stable clip-level supervision, while the selected branch teaches the model to recover text-relevant temporal evidence from imprecise positives. Importantly, multi-level annotation is not used to define a separate hierarchy-specific prediction task. It simply conditions the evidence-selection process through $(K_\ell, \tau_{\mathrm{sel}}^{(\ell)})$. The role of multiple latent queries is therefore not to enforce explicit sentence parsing, but to provide a lightweight evidence-mining mechanism when the caption expresses a compositional procedure description.

\subsection{Hierarchical Consistency Across Adjacent Levels}
\label{sec:method_consistency}

When fine, mid, and coarse samples from the same video are jointly observed during training, their selected views should remain compatible rather than drifting arbitrarily in the embedding space. We therefore construct mixed-level batches and, when triplet anchoring is enabled, include same-video nested triplets satisfying
\begin{equation}
[t_s^{f}, t_e^{f}] \subseteq [t_s^{m}, t_e^{m}] \subseteq [t_s^{c}, t_e^{c}].
\end{equation}
For the available triplets $\mathcal{T}$ in a batch, we apply an adjacent-level consistency loss
\begin{equation}
\mathcal{L}_{\mathrm{hier}} =
\frac{1}{|\mathcal{T}|}
\sum_{(f,m,c)\in \mathcal{T}}
\Big[
1 - \operatorname{cos}(v_{\mathrm{sel}}^f, v_{\mathrm{sel}}^m) +
1 - \operatorname{cos}(v_{\mathrm{sel}}^m, v_{\mathrm{sel}}^c)
\Big].
\end{equation}
We regularize only adjacent levels, not all pairs, because the goal is not to collapse all temporal scopes into one representation. The goal is simply to prevent incompatible local evidence across nested supervision.

\subsection{Training Objective}
\label{sec:method_objective}

The global branch produces
\begin{equation}
v_{\mathrm{glob}} = \operatorname{norm}(P_v(A_v(E_v(c_{\mathrm{exp}})))).
\end{equation}
Both training branches share the same global text embedding $t$ and differ only in how the expanded clip is aggregated on the video side. Using $t$, we optimize two symmetric CLIP-style contrastive objectives:
\begin{equation}
\mathcal{L}_{\mathrm{glob}} =
\mathcal{L}_{v \rightarrow t}(v_{\mathrm{glob}}, t) +
\mathcal{L}_{t \rightarrow v}(t, v_{\mathrm{glob}}),
\end{equation}
\begin{equation}
\mathcal{L}_{\mathrm{sel}} =
\mathcal{L}_{v \rightarrow t}(v_{\mathrm{sel}}, t) +
\mathcal{L}_{t \rightarrow v}(t, v_{\mathrm{sel}}).
\end{equation}
The final objective is
\begin{equation}
\mathcal{L} =
\mathcal{L}_{\mathrm{glob}}
 + \lambda_{\mathrm{sel}} \mathcal{L}_{\mathrm{sel}}
 + \lambda_{\mathrm{hier}} \mathcal{L}_{\mathrm{hier}}.
\end{equation}

Each term has a distinct role. $\mathcal{L}_{\mathrm{glob}}$ preserves stable clip-level video-text alignment over the expanded training window. $\mathcal{L}_{\mathrm{sel}}$ teaches the model to recover local evidence within that window through text-guided reweighting. $\mathcal{L}_{\mathrm{hier}}$ prevents the selected branch from becoming inconsistent across neighboring semantic granularities.

\subsection{Inference}
\label{sec:method_inference}

At inference time, SurgAlign falls back to standard CLIP-style encoding:
\begin{equation}
v = \operatorname{encode\_image}(c), \qquad
t = \operatorname{encode\_text}(y),
\end{equation}
followed by cosine similarity in the shared embedding space. No expanded window, frame reweighting, or hierarchy-specific matching is required at test time. This separation is intentional: SurgAlign is a method for improving weak supervision during pretraining, not a query-conditioned inference system.
